\documentclass[12pt, a4paper]{article}

\usepackage[utf8]{inputenc}
\usepackage[T1]{fontenc}
\usepackage[french]{babel}
\usepackage{geometry}
\usepackage{xcolor}
\usepackage{mdframed}
\usepackage{parskip}
\usepackage{setspace}
\usepackage{titlesec}
\usepackage{csquotes}
\usepackage{hyperref}

\geometry{margin=3cm, top=3.5cm, bottom=3.5cm}
\setstretch{1.4}

\definecolor{accent}{RGB}{24, 95, 165}
\definecolor{lightbg}{RGB}{245, 244, 240}
\definecolor{reddark}{RGB}{121, 31, 31}
\definecolor{greendark}{RGB}{15, 110, 86}

\newmdenv[
  backgroundcolor=lightbg,
  linecolor=accent,
  linewidth=2pt,
  leftline=true, rightline=false, topline=false, bottomline=false,
  innerleftmargin=14pt, innerrightmargin=14pt,
  innertopmargin=10pt, innerbottommargin=10pt,
  skipabove=12pt, skipbelow=12pt
]{encadre}

\newmdenv[
  backgroundcolor=lightbg,
  linecolor=reddark,
  linewidth=2pt,
  leftline=true, rightline=false, topline=false, bottomline=false,
  innerleftmargin=14pt, innerrightmargin=14pt,
  innertopmargin=10pt, innerbottommargin=10pt,
  skipabove=12pt, skipbelow=12pt
]{alerte}

\newmdenv[
  backgroundcolor=lightbg,
  linecolor=greendark,
  linewidth=2pt,
  leftline=true, rightline=false, topline=false, bottomline=false,
  innerleftmargin=14pt, innerrightmargin=14pt,
  innertopmargin=10pt, innerbottommargin=10pt,
  skipabove=12pt, skipbelow=12pt
]{exemple}

\titleformat{\section}{\large\bfseries}{}{0em}{}[\vspace{-0.3em}\rule{\linewidth}{0.4pt}]
\titleformat{\subsection}{\normalsize\bfseries}{}{0em}{}

\hypersetup{colorlinks=true, linkcolor=accent, urlcolor=accent}

\title{
  \textbf{Les routes qui sauvent des vies}\\[0.6em]
  \large Une description du problème posé\\[0.3em]
  \normalsize Sujet 07 — Hackverse 2026
}
\author{}
\date{}

\begin{document}

\maketitle
\vspace{-1em}
\tableofcontents
\newpage


% ───────────────────────────────────────────────────────
\section{La situation de départ}

Dans une grande ville africaine de deux millions d'habitants, quand
quelqu'un compose le numéro des secours, il faut en moyenne vingt-trois
minutes avant qu'une ambulance ou un camion de pompiers arrive sur place.
Vingt-trois minutes. Pour une crise cardiaque, une noyade, un incendie
qui commence à se propager, c'est souvent trop long.

Ce retard n'est pas dû à un manque de véhicules. Les secours existent.
Ils sont disponibles. Le problème est ailleurs : personne ne sait vraiment
quel chemin emprunter. Les conducteurs font appel à leur connaissance du
terrain, à leur expérience du trafic à telle heure, à leur intuition.
Parfois ça marche. Parfois non.

La ville, elle, ne se laisse pas facilement traverser. Les routes changent
de nature selon les quartiers : certaines sont goudronnées et rapides,
d'autres sont des pistes en latérite qui deviennent des bourbiers dès
qu'il pleut. Le trafic suit des habitudes prévisibles le matin et le soir,
mais il peut se figer en quelques minutes à cause d'un accident, d'un
marché qui déborde sur la chaussée, ou d'un camion renversé en plein
carrefour. La nuit, certains axes deviennent dangereux ou impraticables
et les secours doivent les éviter.

\begin{encadre}
L'objectif est donc le suivant : construire un système capable de dire,
à chaque instant, quel est le meilleur chemin pour un véhicule de secours
donné, vers une destination donnée, dans les conditions exactes du moment.
Et si ce chemin se bloque en cours de route, le système doit en trouver un
autre — immédiatement, sans que le conducteur ait besoin de s'arrêter et
de réfléchir.
\end{encadre}


% ───────────────────────────────────────────────────────
\section{Le piège du sujet}

À première lecture, le problème ressemble à une question classique
d'informatique : trouver le plus court chemin dans un réseau. On connaît
les intersections, on connaît les routes qui les relient, on calcule.
Des algorithmes très efficaces existent pour cela. Une équipe motivée
pourrait en construire un en quelques heures.

Mais le sujet contient un avertissement volontairement dissimulé. Il dit,
en substance : une équipe peut avancer vite, produire beaucoup, et
découvrir trop tard qu'elle a soigneusement résolu le mauvais problème.

\begin{alerte}
\textbf{Le vrai piège.}
Calculer le plus court chemin suppose que l'on dispose d'informations
fiables sur l'état du réseau au moment où l'on calcule. Mais dans la
réalité, les informations arrivent en retard, viennent de sources peu
sûres, ou sont tout simplement incomplètes. Le système qui ne tient pas
compte de cela prend des décisions qui semblent bonnes sur le papier,
mais qui peuvent s'avérer désastreuses sur le terrain.
\end{alerte}

Ce n'est donc pas seulement un problème de navigation. C'est un problème
de décision dans l'incertitude. Et ces deux choses sont très différentes.


% ───────────────────────────────────────────────────────
\section{Ce que le système doit vraiment faire}

On peut résumer les attentes en trois mots : choisir, s'adapter,
et expliquer.

\subsection{Choisir}

Le système doit choisir un chemin. Pas le chemin parfait — il n'existe
pas toujours. Mais le meilleur chemin possible avec ce que l'on sait
à cet instant. Ce choix doit respecter des règles absolues : ne jamais
envoyer un véhicule sur un pont qui ne peut pas supporter son poids,
ne jamais engager une ambulance sur une piste inondée la nuit, ne jamais
proposer un itinéraire dont on sait qu'il dépassera le délai maximal
acceptable.

\subsection{S'adapter}

Le monde change pendant que le véhicule roule. Une route se bloque.
La pluie commence à tomber. Un signalement arrive. Le système doit
détecter ces changements et recalculer le chemin sans attendre que
le conducteur se retrouve bloqué. Cette capacité à réagir en temps
réel est l'une des exigences centrales du problème.

\subsection{Expliquer}

C'est ici que le sujet devient original et difficile. Le système doit
être capable de justifier chacune de ses décisions, même après coup.
Si une ambulance arrive en retard, si un chemin s'avère mauvais, si
un responsable demande des comptes un mois plus tard, le système doit
pouvoir répondre : voilà ce que je savais à tel moment, voilà les
informations que j'avais, voilà pourquoi j'ai choisi ce chemin plutôt
qu'un autre.

Ce n'est pas une exigence accessoire. C'est au c\oe{}ur du problème.
Un système qui prend de bonnes décisions mais ne peut pas les justifier
est un système en qui personne ne peut avoir confiance.


% ───────────────────────────────────────────────────────
\section{Une situation concrète : l'ambulance de Bépanda}

Pour rendre tout cela tangible, voici une situation qui illustre les
trois exigences — et qui montre précisément ce qui se passe quand les
informations sont imparfaites.

\subsection{Le départ — tout va bien}

Il est 14h22. Un enfant de sept ans fait des convulsions dans le quartier
Bépanda. L'appel arrive au SAMU. Le système calcule immédiatement le
meilleur chemin : Avenue Bépanda, puis Boulevard de la Liberté, puis le
CHU. Durée estimée : neuf minutes. L'ambulance part. La situation est
nominale.

\subsection{Première situation imparfaite — une information qui arrive trop tard}

Il est 14h25. Le système reçoit une alerte : le Boulevard de la Liberté
est entièrement bloqué par un camion-remorque renversé. Problème : cette
alerte a été générée à 14h21. Elle a mis quatre minutes à traverser le
réseau de collecte surchargé avant d'arriver. Pendant ces quatre minutes,
l'ambulance a continué à rouler vers ce boulevard.

Le système se retrouve donc dans une situation inconfortable. Il sait que
la route est bloquée, mais avec une information vieille de quatre minutes.
Il ne sait pas exactement où se trouve l'ambulance en ce moment — son
signal GPS rame de trente secondes. Il ne sait pas si elle a déjà
franchi le carrefour critique ou si elle peut encore dévier.

Il doit pourtant décider immédiatement. Il évalue trois options : continuer
vers le boulevard en espérant que le blocage soit partiel, prendre la Rue
Makepe qui est une piste non goudronnée mais plus courte que le demi-tour,
ou faire demi-tour complet vers l'Avenue de Gaulle. Le demi-tour complet
dépasse le délai maximum. Le système choisit la Rue Makepe et envoie
l'instruction au conducteur.

\begin{exemple}
\textbf{Ce que le système enregistre à ce moment.}
Il note l'heure exacte de sa décision, l'heure à laquelle l'alerte
lui est parvenue, l'heure à laquelle l'alerte a été générée sur le terrain,
la dernière position connue du véhicule, et les trois options qu'il a
évaluées avec leurs coûts respectifs. Si quelqu'un interroge le système
plus tard, il pourra reconstituer exactement pourquoi ce choix a été fait
avec les informations disponibles à cet instant.
\end{exemple}

\subsection{Deuxième situation imparfaite — une source peu fiable}

Il est 14h27. L'ambulance roule sur la Rue Makepe depuis une minute et
demie. Une nouvelle alerte arrive : un habitant signale via une application
mobile de l'eau sur la chaussée, au niveau du marché, à environ six cents
mètres devant le véhicule. L'alerte ne précise pas la hauteur de l'eau.
Et surtout, ce même signaleur avait envoyé une fausse alerte sept jours
plus tôt, dans le même secteur.

Le conducteur, lui, confirme par radio qu'il voit de l'eau au loin mais
ne peut pas encore juger si c'est passable.

Le système est maintenant dans une situation encore plus délicate que la
première. Il ne peut pas ignorer l'alerte — si l'eau est profonde et que
l'ambulance s'y engage, elle sera bloquée avec l'enfant à bord, loin de
tout. Il ne peut pas non plus y croire aveuglément et ordonner un demi-tour
immédiat — la source est peu fiable, et tout autre itinéraire de substitution
dépasserait le délai maximum.

Il prend alors une décision de compromis : il maintient la route actuelle
de façon provisoire, envoie un message au conducteur lui demandant de
confirmer si le passage est praticable dans les quatre-vingt-dix prochaines
secondes, et recalcule en parallèle un itinéraire de repli prêt à être
envoyé si la confirmation n'arrive pas.

\begin{exemple}
\textbf{Ce que le système enregistre à ce moment.}
Il note la fiabilité estimée de la source, le fait que cette même source
avait déjà produit une fausse alerte, le délai de réponse demandé au
conducteur, et l'itinéraire de repli prêt à l'envoi. La limite de
quatre-vingt-dix secondes est documentée : elle est calculée pour laisser
au conducteur le temps d'évaluer la situation sans compromettre l'arrivée
dans les délais si le repli devient nécessaire.
\end{exemple}

\subsection{La résolution}

À 14h28, soit cinquante-quatre secondes après la demande de confirmation,
le conducteur répond : l'eau mesure environ dix centimètres, le passage
est possible. Le système valide l'itinéraire, annule le repli, et note
que la source du signalement a produit une information exagérée — ce qui
sera pris en compte pour pondérer ses prochaines alertes.

L'ambulance arrive au CHU à 14h31. Neuf minutes après le départ.
L'enfant est pris en charge dans les délais.

\subsection{Ce que ce scénario révèle}

Ce qui rend cette situation instructive, ce n'est pas qu'elle se soit
bien terminée. C'est que le système a dû prendre deux décisions
consécutives sur la base d'informations dégradées — une information en
retard, puis une information de source peu fiable — sans jamais avoir
la certitude que son choix était le bon.

Et pourtant, à chaque étape, sa décision était justifiable : elle était
la meilleure possible avec ce que le système savait à ce moment-là.
C'est exactement ce que le problème demande.


% ───────────────────────────────────────────────────────
\section{Le problème en une phrase}

\begin{encadre}
Construire un moteur capable de guider un véhicule de secours à travers
une ville africaine dont le réseau change en permanence, en prenant les
meilleures décisions possibles avec des informations qui peuvent être
incomplètes, en retard ou peu fiables — et en étant toujours en mesure
d'expliquer, après coup, pourquoi chaque décision a été prise.
\end{encadre}

\end{document}